\section{Reduced-Universe Shell-Rank Selector}
\label{sec:reduced-universe-shell-rank-selector}

The preceding audits isolate the remaining shared\_B selection problem to a finite reduced universe.  The result is deliberately bounded.  It does not derive the full role-labeled shared\_B edge universe from first principles.  It proves an exact selector after the universe has already been reduced to the four observed C-cycles crossed with the four observed twisted anchor-cycles.

\paragraph{Reduced-universe theorem.}
In the 16-candidate shared\_B reduced universe formed by crossing the four observed C-cycles with the four twisted anchor-cycles, define the C shell by the native C-branch marker \(C=5\).  Define the anchor shell by the native branch anchor \([8,18]\).  Rank C-cycles inside each shell by ascending C entry, equivalently by the \(XY.from\_C\) value.  Rank anchor-cycles inside each shell by ascending \(anchor\_sum\_mod15\).  Then the selector
\[
    \mathrm{same\ shell} \;\wedge\; \mathrm{same\ rank}
\]
accepts exactly the four observed shared\_B cycles.

\paragraph{Verification.}
The theorem is verified by artifact 028.  The reduced universe has 16 candidates arranged as a \(4 \times 4\) matrix.  The actual selected candidates are
\[
    [1,6,8,15].
\]
The native shell-rank selector predicts
\[
    [1,6,8,15],
\]
with no false positives and no misses.

The five hypotheses checked in the artifact are:

\begin{enumerate}
\item the reduced universe is a \(4 \times 4\) C-cycle by anchor-cycle universe;
\item the native shell labels are derived;
\item the C-rank rule is unambiguous;
\item the anchor-rank rule is unambiguous;
\item the 027 target shell-rank selector is exact.
\end{enumerate}

All five checks pass.

\paragraph{Proof sketch.}
Artifact 026 derives the shell labels.  On the C side, the branch marker is the repeated XY-to-IW C-junction \(C=5\), giving ordinary C rows \([0,1]\) and branch C rows \([2,3]\).  On the anchor side, the branch marker is the repeated column-motion branch anchor \([8,18]\), giving branch anchor columns \([0,3]\) and ordinary anchor columns \([1,2]\).

Artifact 027 derives the ranks inside those shells.  C rows are ranked by ascending C entry:
\[
    ordinary: 0 < 1,\qquad branch: 2 < 3.
\]
Anchor columns are ranked by ascending \(anchor\_sum\_mod15\):
\[
    branch: 0 < 3,\qquad ordinary: 1 < 2.
\]
Combining the native shell labels with these native ranks gives the four accepted cells
\[
    (0,1),\ (1,2),\ (2,0),\ (3,3),
\]
which are exactly candidates \(1,6,8,15\).

\paragraph{Interpretation.}
The selector is not a scalar threshold.  It is a radial phase-lock: C-cycle shell phase must match anchor-cycle shell phase.  In the matrix representation the selected permutation is
\[
    [1,2,0,3],
\]
with cycle structure \((0\ 1\ 2)(3)\) and order 3.  This matches the geometric reading of a three-phase rotation with one fixed branch shell.



\paragraph{Audit boundary.}
For dependency-audit purposes, we state the boundary in exact terms: this theorem is conditional on the reduced 16-candidate universe. It proves the selector inside the reduced universe. It does not derive the full role-labeled shared\_B edge universe and it does not close Gap A.

\paragraph{Boundary.}
This result does not close Gap A.  The theorem proves the selector inside the reduced 16-candidate universe.  The upstream problem remains: derive the full role-labeled shared\_B edge universe, or equivalently derive the reduced 16-candidate universe itself, from native provenance rather than from observed C-cycle and anchor-cycle support.
